\chapter{Parametrische Modellierung}

\begin{lernziele}
Nach diesem Kapitel verstehen Sie, wie technische Zusammenhänge im Modell beschrieben werden können.
\end{lernziele}

\section{Zweck parametrischer Modelle}
Parametrische Modelle beschreiben mathematische Beziehungen. Sie helfen, technische Entscheidungen zu prüfen. Beim Roboter sind Energieverbrauch, Reichweite und Geschwindigkeit wichtige Größen.

\section{Einfaches Beispiel}
\[
t_{betrieb} = \frac{E_{akku}}{P_{gesamt}}
\]

Dabei ist \(E_{akku}\) die nutzbare Energie der Batterie und \(P_{gesamt}\) die durchschnittliche Leistungsaufnahme.

\section{Leistungsaufnahme}
Die Gesamtleistung kann grob abgeschätzt werden:
\[
P_{gesamt} = P_{rechner} + P_{sensoren} + P_{motoren}
\]

Diese Formel ist einfach, aber nützlich für frühe Architekturentscheidungen.

\section{Parametric Diagram in SysML}
In SysML werden solche Zusammenhänge mit Constraint Blocks modelliert. Ein Constraint Block beschreibt die Formel, Parameter werden mit Systemeigenschaften verbunden.

\begin{hinweisbox}
Parametrik ist für Einsteiger optional. Sie wird wichtig, wenn quantitative Aussagen getroffen werden sollen.
\end{hinweisbox}

\begin{uebungbox}
Nehmen Sie eine Batteriekapazität von 100 Wh und eine Leistungsaufnahme von 25 W an. Berechnen Sie die theoretische Betriebszeit.
\end{uebungbox}
